\documentclass[a4paper,11pt]{article}
\usepackage[utf8]{inputenc}
\usepackage[T1]{fontenc}
\usepackage[ngerman]{babel}
\usepackage{amsmath,amssymb}
\usepackage{xcolor}
\usepackage{colortbl}
\usepackage{array}
\usepackage{tikz}
\usepackage{enumitem}
\usepackage{geometry}
\geometry{a4paper,margin=2cm}
\definecolor{bgdark}{HTML}{1E1E1E}
\definecolor{fgtext}{HTML}{E6E6E6}
\definecolor{accent}{HTML}{6CB6FF}
\definecolor{accent2}{HTML}{F2A65A}
\definecolor{gridcol}{HTML}{4A4A4A}
\definecolor{linecol}{HTML}{8A8A8A}
\pagecolor{bgdark}
\color{fgtext}
\arrayrulecolor{linecol}
\newcommand{\karo}[1]{\par\noindent\begin{tikzpicture}\draw[step=5mm, gridcol, very thin] (0,0) grid (\linewidth,#1);\end{tikzpicture}\par}
\newcommand{\aufgabe}[2]{\medskip\par\noindent{\color{accent}\large\textbf{Aufgabe #1}}\hfill{\color{accent2}\small #2}\par\vspace{0.3em}\hrule height0.4pt\vspace{0.6em}}
\newcommand{\kopf}[1]{\noindent{\color{accent}\Large\textbf{Numerik — #1}}\par\vspace{0.4em}\noindent\textbf{Name:}~\rule[-2pt]{5cm}{0.4pt}\hfill\textbf{Datum:}~\rule[-2pt]{3cm}{0.4pt}\par\vspace{0.3em}\hrule height0.8pt\vspace{1em}}
\setlength{\parindent}{0pt}
\begin{document}
\kopf{Differenzialgleichungen}

Gegeben sei stets eine gewöhnliche Differenzialgleichung erster Ordnung in der
Form $y'=f(y,t)$ mit Anfangswert $y(t_0)=y_0$. Zur Erinnerung:

\textbf{Euler-Verfahren:} \quad $y_{n+1}=y_n+h\cdot f(y_n,t_n)$, \quad $t_{n+1}=t_n+h$.

\textbf{Runge-Kutta 4 (RK4):} \quad $y_{n+1}=y_n+\dfrac h6\,(k_1+2k_2+2k_3+k_4)$, \quad $t_{n+1}=t_n+h$, \quad mit
\[
k_1=f(y_n,t_n),\quad
k_2=f\!\Big(y_n+\tfrac{hk_1}{2},\,t_n+\tfrac h2\Big),\quad
k_3=f\!\Big(y_n+\tfrac{hk_2}{2},\,t_n+\tfrac h2\Big),\quad
k_4=f\!\big(y_n+hk_3,\,t_n+h\big).
\]

\aufgabe{1}{leicht}
Gegeben sei $y'=-\dfrac{5}{y}$ mit Anfangswert $y(0)=5$ und Schrittweite $h=1$.
Berechnen Sie mit dem Euler-Verfahren die Näherungen $y(t_1)=y_1$ und $y(t_2)=y_2$,
also zwei Schritte. Geben Sie auch $t_1$ und $t_2$ an.
\karo{6cm}

\aufgabe{2}{mittel}
Betrachten Sie $y'=-2y+1$ mit $y(0)=0$. Die analytische Lösung lautet
$y(t)=\tfrac12-\tfrac12 e^{-2t}$.
\begin{enumerate}[label=(\alph*)]
\item Rechnen Sie mit dem Euler-Verfahren bis $t=1$, einmal mit $h=0{,}5$ (2 Schritte)
      und einmal mit $h=0{,}25$ (4 Schritte). Notieren Sie jeweils den Näherungswert $y(1)$.
\item Vergleichen Sie beide Näherungen mit dem exakten Wert $y(1)$ und geben Sie an,
      welche Schrittweite näher liegt.
\item Erklären Sie kurz die Fehlerquelle des Euler-Verfahrens (Annahme über die Steigung
      innerhalb eines Schritts).
\end{enumerate}
\karo{8cm}

\aufgabe{3}{leicht}
Wieder sei $y'=-\dfrac{5}{y}$. Nennen Sie einen Startwert $y_0$, der für das
Euler-Verfahren (und für die DGL selbst) problematisch ist, und begründen Sie,
warum dieser Wert nicht verwendet werden darf.
\karo{4cm}

\aufgabe{4}{mittel}
Die DGL $y'=-0{,}01\,y+10$ mit $y(0)=0$ besitzt die analytische Lösung
$y(t)=1000-1000\,e^{-0{,}01\,t}$. Rechnen Sie mit dem Euler-Verfahren zwei Schritte
mit Schrittweite $h=10$, also bis $t=20$. Bestimmen Sie für den Endwert
\textbf{(a)} den absoluten Fehler und \textbf{(b)} den relativen Fehler gegenüber
der analytischen Lösung $y(20)$.
\karo{7cm}

\aufgabe{5}{schwer}
Gegeben sei $y'=f(y,t)=\dfrac{y}{t}+t^2$ mit Anfangswert $y(1)=0{,}5$ und Schrittweite $h=1$.
\begin{enumerate}[label=(\alph*)]
\item Berechnen Sie für \emph{einen} RK4-Schritt alle vier Koeffizienten
      $k_1,k_2,k_3,k_4$ von Hand (mit Brüchen).
\item Bestimmen Sie daraus $y_{n+1}=y_n+\dfrac h6(k_1+2k_2+2k_3+k_4)$.
\item Vergleichen Sie mit der exakten Lösung $y(t)=0{,}5\,t^3$ (also $y(2)=4$).
\end{enumerate}
\karo{8cm}

\aufgabe{6}{mittel}
Betrachten Sie $y'=y$ mit $y(0)=1$ und Schrittweite $h=1$. Die exakte Lösung ist
$y(t)=e^{t}$, also $y(1)=e\approx 2{,}71828$.
\begin{enumerate}[label=(\alph*)]
\item Rechnen Sie einen Schritt mit dem Euler-Verfahren.
\item Rechnen Sie einen Schritt mit RK4 (alle $k_i$ angeben).
\item Vergleichen Sie beide Näherungen mit $e$ und erklären Sie, warum RK4 deutlich
      genauer ist (Stichwort: Auswertung der Steigung an mehreren Punkten).
\end{enumerate}
\karo{8cm}

\aufgabe{7}{mittel}
Gegeben sei das folgende Butcher-Tableau:
\[
\renewcommand{\arraystretch}{1.3}
\begin{array}{c|cc}
0 & & \\
\tfrac23 & \tfrac23 & \\
\hline
 & \tfrac14 & \tfrac34
\end{array}
\]
\begin{enumerate}[label=(\alph*)]
\item Wie viele Stufen hat das Verfahren?
\item Lesen Sie die Formeln für $k_1$ und $k_2$ sowie die Aktualisierungsformel
      $y_{n+1}$ aus dem Tableau ab.
\item Welche Bedingung an die Gewichte muss erfüllt sein? Prüfen Sie sie.
\end{enumerate}
\karo{7cm}

\aufgabe{8}{mittel}
Stellen Sie das Butcher-Tableau für das klassische RK4-Verfahren auf. Erklären Sie
dabei die Bedeutung der drei Bestandteile (Knoten in der linken Spalte, Matrix rechts
oben, Gewichtszeile unten). Begründen Sie, warum die Gewichte
$\tfrac16,\tfrac26,\tfrac26,\tfrac16$ lauten und ihre Summe gleich $1$ sein muss.
\karo{8cm}

\aufgabe{9}{schwer}
\begin{enumerate}[label=(\alph*)]
\item Zeigen Sie ausgehend von der RK4-Formel, dass man genau das Euler-Verfahren
      erhält, wenn man $k_2,k_3,k_4$ weglässt und nur $k_1$ mit Vorfaktor $h$ behält.
\item Diskutieren Sie den Zusammenhang RK1 $=$ Euler: Welche „Stufe`` besitzt das
      Euler-Verfahren und wie sieht sein Butcher-Tableau aus?
\item (Optional) Geben Sie das Butcher-Tableau der Mittelpunktregel (RK2) an und
      erläutern Sie, warum sie ein Spezialfall der Runge-Kutta-Familie ist.
\end{enumerate}
\karo{8cm}

\end{document}
